% Created 2026-05-12 ter 18:31
% Intended LaTeX compiler: lualatex
\documentclass[12pt,a4paper]{article}
\usepackage{fgv-paper}
% Bibliografia controlada pelo TOML; os estilos FGV não carregam biblatex.
\usepackage[backend=biber,sorting=nty,giveninits=true]{biblatex}

\addbibresource{resumo_artigos_encontro_2.bib}
\geometry{a4paper,left=3cm,right=2cm,top=3cm,bottom=2cm,headheight=58pt,headsep=18pt,footskip=1.2cm}
\usepackage{float}
\usepackage{tikz}
\usepackage{tcolorbox}
\usepackage{tabularray}
\usepackage{ltablex}
\keepXColumns
\setlength{\emergencystretch}{3em}
\sloppy
\definecolor{FGVHeaderBlueDark}{HTML}{003B71}
\definecolor{FGVHeaderBlueMid}{HTML}{006BB6}
\definecolor{FGVHeaderBlueLight}{HTML}{8FD3F4}
\setlength{\headheight}{58pt}
\setlength{\headsep}{18pt}
\newcommand{\FGVHeaderLogoInclude}{\IfFileExists{/home/gustavodetarso/Documentos/mppg/software/academic_pipeline_rc10_7_conformidade/app_bundle/institutions/fgv/assets/fgv.png}{
\includegraphics[height=1.05cm]{/home/gustavodetarso/Documentos/mppg/software/academic_pipeline_rc10_7_conformidade/app_bundle/institutions/fgv/assets/fgv.png}}{}}
\newcommand{\FGVHeaderGradientRule}{\begin{tikzpicture}[baseline=(current bounding box.center)]\shade[left color=FGVHeaderBlueDark,middle color=FGVHeaderBlueMid,right color=FGVHeaderBlueLight] (0pt,0pt) rectangle (\headwidth,2pt);\end{tikzpicture}}
\newcommand{\FGVHeaderBlock}{\makebox[\headwidth][l]{\begin{minipage}[t]{\headwidth}\FGVHeaderLogoInclude\\[-1pt]\FGVHeaderGradientRule\end{minipage}}}
\fancypagestyle{fgvreportstyle}{%
\fancyhf{}%
\fancyhead[L]{\FGVHeaderBlock}%
\renewcommand{\headrulewidth}{0pt}%
\fancyfoot[C]{\thepage}%
}
\fancypagestyle{plain}{%
\fancyhf{}%
\fancyhead[L]{\FGVHeaderBlock}%
\renewcommand{\headrulewidth}{0pt}%
\fancyfoot[C]{\thepage}%
}
\AtBeginDocument{\pagestyle{fgvreportstyle}\thispagestyle{fgvreportstyle}}
\institution{fgv}
\programname{}
\coursename{Mestrado Acadêmico em Políticas Públicas e Governo}
\disciplinename{Políticas Públicas - Teorias e Análise}
\professorname{Lizandro Lui}
\cityname{Brasília}
\author{Gustavo M. Mendes de Tarso}
\date{\today}
\title{Resumo analítico dos textos}
\hypersetup{
 pdfauthor={Gustavo M. Mendes de Tarso},
 pdftitle={Resumo analítico dos textos},
 pdfkeywords={},
 pdfsubject={},
 pdfcreator={Emacs 28.2 (Org mode 9.5.5)}, 
 pdflang={Pt_Br}}
\begin{document}

\begingroup
\linespread{1}\selectfont
\begin{tcolorbox}[title=Ficha Técnica,
  colback=gray!5,colframe=gray!40,boxrule=0.4pt,sharp corners]
\begin{tblr}{rowsep=1pt,stretch=1, rows={t}, colspec={Q[l,2.8cm] X[l]}}
\textbf{Curso}                   & Mestrado Acadêmico em Políticas Públicas e Governo \\
\textbf{Turma}                   & 2026.1 \\
\textbf{Pólo}                    & Brasília \\
\textbf{Disciplina}              & Políticas Públicas - Teorias e Análise \\
\textbf{Professor}               & Lizandro Lui \\
\textbf{Aluno(s)}                & Gustavo M. Mendes de Tarso \\
\textbf{Data}                    & 2026 \\
\textbf{Título do trabalho}      & Resumo analítico dos textos \\
\end{tblr}
\end{tcolorbox}
\endgroup

\clearpage

\section{Introdução}
\label{sec:org4806d2a}

O corpus examinado é composto por dois textos que se situam no campo das políticas públicas, mas operam em níveis analíticos distintos e, por isso mesmo, oferecem uma combinação intelectualmente produtiva. Um deles reconstrói e problematiza os principais modelos de análise da implementação, propondo ampliar o repertório explicativo do campo; o outro acompanha a evolução de uma política pública específica, voltada ao enfrentamento do desmatamento na Amazônia Legal, com atenção às fases de sua implementação, aos instrumentos mobilizados e aos resultados alcançados. \parencite{ascenzi2013_implementacao_politicas,mello2017_evolucao_plano}

O foco de leitura adotado neste documento recai sobre a implementação como problema analítico, e não apenas como etapa administrativa posterior à formulação. Isso significa observar como cada texto, a seu modo, ajuda a compreender de que depende a efetivação de uma política: se da clareza do desenho institucional, da coordenação entre órgãos, das capacidades estatais mobilizadas, das condições locais de execução ou também das interpretações e valores dos atores envolvidos. Em vez de tratar os textos como peças isoladas, a análise procura extrair deles um diálogo sobre os modos de explicar a ação pública em contextos complexos. \parencite{ascenzi2013_implementacao_politicas,mello2017_evolucao_plano}

Há, portanto, heterogeneidade real no corpus. Essa heterogeneidade não deve ser confundida com dispersão improdutiva. Os textos não compartilham o mesmo objeto empírico nem a mesma escala de observação: um trabalha em chave teórica e reconstrutiva; o outro, em chave empírica e longitudinal. Ainda assim, ambos convergem em um eixo disciplinar comum: recusam a imagem de implementação como execução mecânica e insistem que políticas públicas se transformam quando entram em contato com organizações, territórios, instrumentos, burocracias e dinâmicas concretas de ação. A unidade do corpus, assim, não é temática no sentido estreito, mas analítica no sentido do problema que mobiliza a leitura. \parencite{ascenzi2013_implementacao_politicas,mello2017_evolucao_plano}

O objetivo deste documento é produzir uma leitura analítica aprofundada do conjunto, fiel aos textos e ao mesmo tempo interpretativa. Não se trata apenas de resumir argumentos, mas de reconstruir o que cada trabalho efetivamente demonstra, identificar sua contribuição específica para a disciplina e, sobretudo, explicitar como eles podem ser articulados sem apagamento de suas diferenças. Com isso, busca-se oferecer uma síntese útil para o estudo de políticas públicas, especialmente no que diz respeito à relação entre formulação, implementação, capacidades estatais, contexto e mediações interpretativas. \parencite{ascenzi2013_implementacao_politicas,mello2017_evolucao_plano}

O caminho analítico adotado parte, primeiro, do reconhecimento da especificidade de cada texto: um como intervenção conceitual sobre o estudo da implementação, outro como reconstrução de uma trajetória concreta de política ambiental. Em seguida, toma-se a comparação apenas nos eixos efetivamente sustentados pelo corpus — problema abordado, argumento central, conceitos mobilizados, tipo de abordagem, método, contribuição e limites. Por fim, a síntese interpreta o conjunto a partir de uma hipótese de leitura: a de que a implementação só pode ser compreendida de forma robusta quando capacidades institucionais e mediações interpretativas são pensadas em articulação, ainda que cada texto enfatize dimensões distintas dessa relação. \parencite{ascenzi2013_implementacao_politicas,mello2017_evolucao_plano}



\section{Resumo individual dos textos}
\label{sec:orgc26568b}



\subsection{Implementação de políticas públicas: perspectivas analíticas}
\label{sec:orge58c83d}

O artigo enfrenta uma questão central do campo de políticas públicas: como explicar a trajetória concreta da implementação sem reduzi-la nem a mera execução automática de decisões previamente formuladas, nem a um processo inteiramente contingente e indeterminado. A intervenção dos autores parte do diagnóstico de que a literatura clássica organizou o debate em torno do “problema da implementação”, isto é, da busca por variáveis capazes de explicar sucesso ou fracasso na realização dos objetivos estatais. Em vez de simplesmente reiterar esse enquadramento, o texto procura mostrar que a própria maneira de definir o problema orienta o recorte analítico, os fatores considerados relevantes e as soluções propostas. A contribuição inicial do artigo é, portanto, reflexiva: antes de escolher variáveis, é preciso interrogar os pressupostos embutidos nos modelos de análise da implementação\parencite{ascenzi2013_implementacao_politicas}

O objetivo do texto é duplo e explicitamente articulado. Primeiro, sistematizar e analisar criticamente a literatura sobre modelos analíticos de implementação de políticas públicas, identificando seus núcleos explicativos dominantes. Segundo, propor uma abordagem integradora que preserve a utilidade das variáveis clássicas, mas incorpore dimensões cognitivas e ideacionais, especialmente ideias, valores, representações e concepções de mundo dos atores envolvidos na execução. O escopo do artigo é, assim, ao mesmo tempo reconstrutivo e propositivo: reconstrói o debate para evidenciar suas clivagens principais e, a partir de pesquisas empíricas dos próprios autores em saúde e educação, sugere um alargamento do repertório analítico disponível. O texto não se limita a revisar autores; ele procura reposicionar o problema da implementação em bases menos restritivas e mais sensíveis aos contextos locais de ação\parencite{ascenzi2013_implementacao_politicas}

A tese principal sustenta que os dois modelos hegemônicos da literatura, embora úteis, compartilham uma limitação importante: ambos enfatizam sobretudo condições materiais e organizacionais do processo implementador e deixam em segundo plano a mediação simbólica pela qual os atores interpretam problemas, objetivos e meios de ação. O primeiro modelo privilegia a hierarquia organizacional e o desenho da política, concentrando-se nas normas, na clareza de objetivos, na definição de responsabilidades e no controle da discricionariedade. O segundo desloca a atenção para os contextos locais e para as burocracias implementadoras, reconhecendo maior autonomia dos executores e maior peso das interações situadas. A inovação do artigo está em argumentar que, mesmo nesse segundo registro, ainda falta incorporar de modo sistemático as representações e valores que estruturam a percepção dos agentes sobre o problema social e sobre as formas legítimas ou adequadas de enfrentá-lo\parencite{ascenzi2013_implementacao_politicas}

Na reconstrução do primeiro polo analítico, o artigo descreve a perspectiva orientada pelo processo de formulação, frequentemente associada ao modelo top-down. Nessa visão, a implementação é concebida como decorrência da decisão anterior e, portanto, sua explicação remete prioritariamente à qualidade do desenho da política. A boa implementação dependeria de objetivos claros, teoria causal consistente, estrutura normativa precisa, recursos adequados, apoio político e aparato administrativo capaz de transmitir e fazer cumprir comandos. O parâmetro de sucesso é o grau de aderência entre resultados observados e metas formuladas. Ao reconstruir esse modelo, os autores não o descartam sumariamente; reconhecem sua relevância para pensar problemas de coordenação, desenho institucional e capacidade estatal. Contudo, mostram que ele repousa sobre um pressuposto forte e problemático: a ideia de que formuladores controlam, ou ao menos podem estruturar de forma decisiva, os processos organizacionais e políticos que condicionam a implementação\parencite{ascenzi2013_implementacao_politicas}

O segundo polo enfatizado no artigo desloca a análise para os espaços locais de implementação e para a atuação das burocracias encarregadas da execução. Nesse registro, a implementação deixa de ser vista como etapa subordinada e tecnicamente neutra, passando a ser compreendida como processo de interação, negociação, adaptação e, em certos casos, redefinição da própria política. A ambiguidade de objetivos, a limitação de recursos, os problemas de coordenação e a pluralidade de interesses tornam improvável a imagem de um fluxo linear entre formulação e execução. O mérito desse modelo está em reconhecer o peso da discricionariedade e a centralidade da burocracia implementadora, sobretudo em contextos nos quais o contato com públicos, territórios e organizações concretas produz ajustes contínuos. Ainda assim, os autores sugerem que a ênfase no contexto e na ação das burocracias, embora avance em relação ao top-down, não esgota o problema explicativo, pois carece de um tratamento mais refinado sobre como os atores interpretam a política que executam\parencite{ascenzi2013_implementacao_politicas}

Do ponto de vista conceitual, o texto mobiliza categorias decisivas para a disciplina. Entre elas, destacam-se implementação, discricionariedade, burocracia implementadora, contexto local, variáveis independentes, sucesso e fracasso da política, além da oposição entre desenho prospectivo e abordagens sensíveis à dinâmica da execução. Mas o movimento mais importante do artigo é introduzir, como categoria analítica relevante, a dimensão cognitiva e ideológica da implementação. Isso significa afirmar que valores, ideias e concepções de mundo não são resíduos subjetivos ou ruídos externos ao processo, mas elementos constitutivos da prática administrativa. Ao interpretarem o problema social, os implementadores também interpretam a política, seus destinatários e os meios considerados apropriados para agir. Com isso, a implementação aparece menos como tradução mecânica de comandos e mais como atividade situada de construção de sentido, ainda que sempre condicionada por normas, recursos e estruturas organizacionais\parencite{ascenzi2013_implementacao_politicas}

Em termos de método e estratégia argumentativa, o artigo combina revisão analítica da literatura com apoio em pesquisas empíricas anteriores dos próprios autores sobre políticas de saúde e educação. Não se trata de um estudo de caso único nem de teste formal de hipóteses, mas de uma operação teórico-analítica sustentada por evidências empíricas mobilizadas para demonstrar os limites dos modelos tradicionais e fundamentar a proposta de integração. Essa estratégia é coerente com o propósito do texto: mais do que medir efeitos isolados, interessa construir um modelo interpretativo capaz de orientar futuras investigações. As pesquisas empíricas mencionadas funcionam como prova de plausibilidade da hipótese de que representações e valores dos atores, nos contextos locais de ação, interferem na execução das políticas ao moldarem a compreensão do problema e da resposta cabível. O artigo, portanto, produz generalização analítica por meio de síntese teórica informada pela experiência de pesquisa\parencite{ascenzi2013_implementacao_politicas}

Os achados e contribuições do texto concentram-se precisamente nessa ampliação do horizonte explicativo. Ao mostrar que contextos locais não são apenas arranjos materiais ou institucionais, mas também universos de significação, o artigo desloca o foco da implementação para a interface entre estrutura e interpretação. A principal contribuição para a disciplina é oferecer uma proposta de integração entre modelos frequentemente tratados como alternativos ou excludentes. Em vez de escolher entre desenho normativo e dinâmica local, os autores defendem uma composição analítica em que variáveis organizacionais, recursos, regras e capacidades importam, mas sua eficácia depende também dos esquemas cognitivos pelos quais os atores compreendem o que está em jogo. Para os estudos de administração pública e políticas públicas, isso implica reconhecer que a ação estatal se realiza por meio de agentes que não apenas obedecem ou resistem, mas atribuem sentido ao seu trabalho e ao público com o qual interagem\parencite{ascenzi2013_implementacao_politicas}

Essa contribuição também produz implicações importantes para a operacionalização da pesquisa em implementação. Se a compreensão dos atores sobre o problema social afeta a execução da política, então análises centradas apenas em desenho institucional, recursos e cadeias hierárquicas tendem a capturar apenas parte do processo. O artigo sugere, assim, a necessidade de observar os espaços locais de ação de maneira mais densa, investigando como burocratas interpretam normas, classificam situações, percebem destinatários e legitimam determinadas soluções. A disciplina ganha, com isso, uma agenda mais fina para estudos empíricos: a implementação deve ser analisada como processo no qual regras e ideias se entrelaçam. Em termos práticos, a proposta ajuda a compreender por que políticas formalmente semelhantes podem produzir trajetórias distintas quando executadas em ambientes organizacionais e simbólicos diferentes, mesmo sob desenhos normativos relativamente estáveis\parencite{ascenzi2013_implementacao_politicas}

Há, contudo, limites e tensões relevantes no próprio artigo. A defesa da integração entre variáveis materiais, organizacionais e cognitivas é intelectualmente fecunda, mas permanece em nível relativamente programático. O texto convence mais quanto à insuficiência dos modelos tradicionais do que quanto à plena especificação do modelo alternativo. Fica clara a pertinência de incluir ideias e valores, porém menos definidos os mecanismos pelos quais esses elementos se articulam com estruturas institucionais, recursos e incentivos, bem como os procedimentos para sua mensuração ou comparação sistemática. Além disso, como a base empírica aparece de forma ilustrativa e não extensivamente apresentada no artigo, a proposta depende de desenvolvimentos posteriores para ganhar maior precisão metodológica. Esses limites, entretanto, não reduzem sua relevância; ao contrário, indicam que a principal força do texto está em abrir uma agenda de pesquisa e não em encerrá-la\parencite{ascenzi2013_implementacao_politicas}

Outro ponto de tensão decorre do próprio vocabulário de sucesso e fracasso que organiza parte do debate revisado. Embora o artigo se distancie de leituras estritamente normativas da implementação, ele ainda dialoga com um campo fortemente orientado por expectativas de controle, melhoria e correção de problemas. Sua inflexão cognitiva relativiza essa lógica ao mostrar que a implementação transforma a política e é transformada por interpretações situadas, mas não abandona completamente o enquadramento de desempenho. Isso produz uma ambivalência produtiva: de um lado, preserva a utilidade aplicada do campo; de outro, aproxima a análise de uma sociologia da ação pública mais atenta à construção social dos problemas e das respostas. Para a disciplina, essa ambivalência é menos uma falha do que um indício da posição de fronteira ocupada pelo artigo entre análise orientada à gestão e interpretação sociopolítica dos processos estatais\parencite{ascenzi2013_implementacao_politicas}

No diálogo com o restante do corpus, o texto oferece uma moldura analítica especialmente útil para ler estudos aplicados sobre trajetórias concretas de políticas públicas. Enquanto o outro trabalho do conjunto se concentra na evolução da implementação de um plano específico, examinando fases, instrumentos e rearranjos institucionais ao longo do tempo, este artigo opera em nível mais conceitual e metateórico, fornecendo categorias para interpretar justamente por que políticas mudam em sua execução e por que seus resultados dependem de múltiplas mediações\parencite{ascenzi2013_implementacao_politicas}; ver também\parencite{mello2017_evolucao_plano}

A relação entre os dois textos, porém, não deve ser forçada como se ambos compartilhassem a mesma escala ou a mesma ambição teórica. O artigo aqui analisado é uma intervenção de sistematização e proposição conceitual; o outro se orienta por um objeto substantivo e por uma reconstrução histórica da implementação de uma política ambiental. A aproximação real entre eles está no reconhecimento de que implementação não é simples fase posterior à formulação, mas processo atravessado por contextos, capacidades estatais, coordenação e variações ao longo do tempo. Se este texto insiste em incorporar ideias e valores dos atores como parte do modelo explicativo, ele oferece uma lente promissora para aprofundar análises empíricas de políticas setoriais, inclusive aquelas em que a variação da implementação decorre não apenas de desenho institucional, mas também de interpretações concorrentes sobre o problema público e sobre a intervenção estatal legítima\parencite{ascenzi2013_implementacao_politicas,mello2017_evolucao_plano}



\subsection{Evolução do Plano de Ação para Prevenção e Controle do Desmatamento na Amazônia Legal}
\label{sec:orgf097be2}

O artigo examina a trajetória do Plano de Ação para Prevenção e Controle do Desmatamento na Amazônia Legal como resposta estatal a um problema que, no início dos anos 2000, alcançou o centro da agenda governamental brasileira. Seu ponto de partida é menos a descrição episódica do plano do que a reconstrução do processo pelo qual o desmatamento deixa de ser apenas uma condição social e ambiental preocupante para se consolidar como problema público que demanda coordenação governamental de alta densidade institucional. Nesse enquadramento, o texto sugere que compreender o PPCDAm exige articular a gravidade ecológica da Amazônia, a elevação das taxas de desmatamento, a pressão política por resposta e a formação de uma arquitetura interministerial capaz de converter diagnóstico em ação pública \parencite{mello2017_evolucao_plano}

O objetivo do estudo é explicitamente analítico e evolutivo: acompanhar a implementação do PPCDAm ao longo de suas três fases, comparar suas diretrizes estratégicas e planos operacionais, e verificar em que medida avaliações independentes influenciaram ajustes subsequentes. O escopo, portanto, combina duas dimensões. De um lado, há uma reconstrução longitudinal da política, que permite observar mudança institucional, redefinição de prioridades e aprendizado administrativo. De outro, há uma preocupação avaliativa voltada aos efeitos do plano sobre a redução do desmatamento e à insuficiência relativa de seus mecanismos para promover atividades sustentáveis na região. O texto não se limita a perguntar se a política funcionou; ele pergunta em que termos funcionou, por quais eixos de ação e com quais déficits persistentes \parencite{mello2017_evolucao_plano}

A tese principal é equilibrada, mas nítida: o PPCDAm produziu resultados significativos na contenção do desmatamento na Amazônia Legal, especialmente por meio de mecanismos de monitoramento, comando e controle, ordenamento territorial e coordenação federativa e intersetorial; contudo, sua evolução revela que a dimensão de promoção de atividades produtivas sustentáveis permaneceu comparativamente mais frágil. Isso significa que o plano é apresentado como caso bem-sucedido de resposta estatal à emergência ambiental, mas não como solução integral da questão amazônica. A contribuição do artigo está justamente em evitar tanto a celebração acrítica da queda do desmatamento quanto a desqualificação simplista da política: o argumento sustenta que houve êxito concreto, embora seletivo e incompleto \parencite{mello2017_evolucao_plano}

Do ponto de vista conceitual, o artigo mobiliza uma linguagem de políticas públicas para explicar a emergência do tema na agenda governamental. Ao recuperar definições de política pública e, sobretudo, a formulação dos múltiplos fluxos, os autores constroem uma moldura analítica em que problema, alternativas e dinâmica política convergem para possibilitar ação estatal. Essa escolha não é meramente introdutória. Ela organiza a interpretação de que a escalada do desmatamento, associada à relevância ecossistêmica da Amazônia e à disponibilidade de instrumentos estatais e diagnósticos prévios, criou uma janela de oportunidade para a institucionalização do plano. O texto, assim, insere o PPCDAm em uma gramática de agenda, formulação e implementação, ainda que seu foco empírico recaia sobretudo sobre esta última dimensão \parencite{mello2017_evolucao_plano}

A caracterização do problema do desmatamento é construída de maneira multifatorial. O texto mostra que as mudanças no uso do solo são espacial e temporalmente heterogêneas e resultam de longos ciclos de ocupação, infraestrutura, incentivos econômicos, colonização, crédito e grilagem. Ao reconstituir esse pano de fundo, os autores afastam leituras monocausais e ressaltam que o desmatamento não decorre apenas de comportamentos individuais ilegais, mas de uma trajetória de desenvolvimento regional marcada por políticas estatais pretéritas e por assimetrias na governança territorial. Esse enquadramento é importante porque reforça a ideia de que a implementação do PPCDAm precisaria ser necessariamente integrada, interministerial e sensível à complexidade causal do problema \parencite{mello2017_evolucao_plano}

O método empregado é de natureza documental e analítica, baseado na comparação entre publicações científicas sobre o problema, documentos oficiais das três fases do plano, avaliações independentes encomendadas pelo poder público e estudos sobre os impactos do PPCDAm na redução do desmatamento. A estratégia argumentativa é cumulativa: primeiro caracteriza-se o problema; em seguida, reconstrói-se a arquitetura e a evolução das fases do plano; por fim, cotejam-se avaliações e resultados para identificar quais eixos foram mais efetivos. Não se trata de um desenho causal estrito nem de avaliação econométrica original, mas de uma síntese analítica orientada por comparação longitudinal de documentos e evidências secundárias. Essa opção dá ao artigo amplitude institucional e histórica, embora limite o grau de identificação causal fina \parencite{mello2017_evolucao_plano}

Na reconstrução das fases do PPCDAm, o texto evidencia um processo de amadurecimento institucional. A primeira fase aparece como momento de estruturação do esforço estatal, com ênfase na articulação interministerial e na organização de ações integradas. A segunda sinaliza refinamento, consolidação de instrumentos e incorporação de aprendizados oriundos das avaliações. A terceira indica tentativa de aprofundar e estabilizar a política, inclusive no horizonte de conciliar controle do desmatamento com promoção de bases sustentáveis para o desenvolvimento regional. O artigo lê essa passagem entre fases como evidência de adaptação institucional: a política não permanece idêntica a si mesma, mas é reelaborada à luz de resultados, críticas e novos desafios \parencite{mello2017_evolucao_plano}

Os achados mais relevantes concentram-se na identificação dos eixos de ação com maior capacidade de produzir efeitos. A análise aponta que o plano foi particularmente eficaz quando combinou monitoramento ambiental, fiscalização, restrições administrativas, presença coordenada do Estado e instrumentos de ordenamento territorial. Essa leitura importa para a disciplina porque desloca a discussão de uma oposição simplista entre formulação e execução para uma questão mais concreta: quais capacidades estatais, em quais arranjos institucionais, conseguem alterar trajetórias de uso do solo em territórios extensos, heterogêneos e atravessados por interesses econômicos conflitantes. Nesse sentido, o artigo oferece um caso substantivo de política pública complexa cuja efetividade depende de coordenação horizontal, informação territorial e capacidade de enforcement \parencite{mello2017_evolucao_plano}

Ao mesmo tempo, o texto deixa claro que a contenção do desmatamento não equivale automaticamente à consolidação de um modelo sustentável para a Amazônia Legal. A dimensão de fomento a atividades econômicas sustentáveis surge como a face menos robusta do plano, o que produz uma tensão central: políticas de controle podem reduzir a velocidade da destruição florestal sem, por si sós, reordenar de maneira duradoura os incentivos econômicos e sociais que estruturam o território. A implicação é relevante porque sugere que o sucesso do PPCDAm, embora real, repousa fortemente sobre instrumentos repressivos e regulatórios, carecendo de complemento mais consistente na esfera do desenvolvimento sustentável. O artigo, portanto, ajuda a distinguir entre redução de dano e transformação estrutural \parencite{mello2017_evolucao_plano}

Entre os limites do trabalho, destaca-se precisamente o fato de que sua síntese depende de documentação oficial e avaliações já produzidas, o que fortalece a reconstrução institucional, mas oferece menor acesso às microdinâmicas de implementação nos territórios, às disputas entre burocracias, aos efeitos diferenciais entre estados e municípios e às respostas dos atores locais às medidas do plano. Há também um desbalanceamento inerente entre a demonstração dos resultados agregados de queda do desmatamento e a exploração mais fina das mediações causais que conectam instrumentos específicos a resultados específicos. Ainda assim, tais limites não anulam a contribuição do artigo; antes, delimitam seu lugar como análise abrangente da evolução de uma política, mais do que como estudo de campo sobre implementação situada \parencite{mello2017_evolucao_plano}

No diálogo com o outro texto do corpus, a aproximação mais produtiva está no tema da implementação. Enquanto o estudo sobre perspectivas analíticas da implementação de políticas públicas opera em nível teórico e procura sistematizar modelos explicativos, propondo integrar variáveis institucionais e materiais com ideias, valores e concepções de mundo dos atores, este artigo oferece um caso empírico de política complexa em que coordenação organizacional, capacidades estatais e contextos territoriais são decisivos para os resultados obtidos. A comparação mostra convergência quanto à insuficiência de leituras lineares da ação pública, mas também revela diferença de ênfase: aqui predominam desenho institucional, fases do plano e instrumentos de ação, ao passo que no texto teórico há maior esforço para explicitar o papel das mediações cognitivas e normativas dos agentes implementadores. Em conjunto, os dois textos não formam uma unidade artificial, mas se iluminam mutuamente: um explicita as lentes de análise da implementação; o outro demonstra como uma política de alta complexidade pode ser reconstruída a partir de seus arranjos, resultados e tensões \parencite{ascenzi2013_implementacao_politicas,mello2017_evolucao_plano}



\section{Comparação entre os textos}
\label{sec:org7de5a67}

\begin{landscape}
\begingroup
\scriptsize
\setlength{\tabcolsep}{3pt}
\renewcommand{\arraystretch}{1.12}
\begin{longtblr}[caption={Quadro comparativo sintético dos textos do corpus},label={tbl-comparacao-corpus}]{%
  width=\linewidth,
  colspec={X[1.0,l] X[1.2,l] X[1.15,l] X[1.15,l] X[1.25,l] X[1.15,l]},
  rowhead=1,
  rows={valign=t},
  row{1}={font=\bfseries},
  hline{1,Z}={0.08em},
  hline{2}={0.04em},
}
Texto & Objeto/problema & Escala e método & Conceito central & Contribuição & Limite/tensão \\
Perspectivas analíticas da implementação & Modelos de análise da implementação de políticas públicas & Escala teórica; revisão analítica com apoio empírico ilustrativo & Integração entre variáveis organizacionais, contextuais e dimensões ideacionais & Reformula o estudo da implementação para além do desenho formal e das condições materiais & Modelo proposto permanece mais programático do que plenamente operacionalizado \\
Evolução do plano de prevenção e controle do desmatamento & Trajetória de implementação de um plano ambiental federal & Escala meso e macro; análise documental longitudinal de fases, avaliações e resultados & Coordenação estatal, monitoramento, comando e controle, ordenamento territorial & Mostra como capacidades estatais e coordenação intersetorial produziram redução substantiva do desmatamento & Explora menos as mediações locais, cognitivas e territoriais da implementação concreta \\
\end{longtblr}
\endgroup
\end{landscape}

Os dois textos convergem no ponto mais importante do corpus: ambos recusam a visão da implementação como etapa mecânica, linear e inteiramente subordinada à formulação. Ainda que o façam em registros distintos, os trabalhos sustentam que a execução de políticas públicas depende de mediações organizacionais, institucionais e contextuais que alteram, redefinem ou ao menos modulam o que a política se torna na prática. No texto de perspectiva teórica, essa crítica aparece por meio da reconstrução dos modelos analíticos da literatura e da proposição de uma abordagem integradora; no estudo sobre o plano voltado ao desmatamento, ela emerge empiricamente pela demonstração de que os resultados da política dependeram de arranjos de coordenação, adaptação entre fases, aprendizado institucional e combinação de instrumentos. Assim, há uma convergência real quanto à rejeição de qualquer leitura simplificada da implementação, embora os caminhos argumentativos sejam bastante diferentes \parencite{ascenzi2013_implementacao_politicas,mello2017_evolucao_plano}

A divergência mais evidente diz respeito ao objeto de análise. Um dos artigos toma a própria literatura de implementação como problema e se dedica a examinar seus modelos explicativos, suas variáveis privilegiadas e suas insuficiências. Seu objeto não é uma política específica, mas a arquitetura conceitual pela qual o campo procura explicar o sucesso, o fracasso e a variação dos processos implementadores. Já o outro texto investiga uma política pública determinada, com recorte substantivo muito claro: a evolução de um plano federal de enfrentamento do desmatamento na Amazônia Legal ao longo de diferentes fases. Em outras palavras, um trabalho se move no plano metateórico da disciplina; o outro, no plano empírico de uma trajetória concreta de ação estatal. Essa diferença de objeto impede comparações apressadas e exige reconhecer que os textos não competem entre si, mas operam em níveis analíticos distintos \parencite{ascenzi2013_implementacao_politicas,mello2017_evolucao_plano}

Essa diferença de objeto se articula a uma diferença igualmente decisiva de escala analítica. O texto de reconstrução teórica trabalha em escala elevada de abstração. Seu interesse está em identificar famílias de explicação, isolar variáveis centrais e demonstrar que tanto abordagens mais centradas na hierarquia e no desenho quanto abordagens voltadas ao contexto local e à burocracia implementadora compartilham um viés material e organizacional. A partir daí, propõe incorporar ideias, valores e concepções de mundo como dimensão constitutiva da implementação. Já o estudo sobre o plano ambiental opera em escala meso e macro, acompanhando uma política nacional complexa, intersetorial e territorialmente extensa, com atenção às suas fases, aos seus instrumentos e aos seus resultados agregados. O primeiro texto pergunta como pensar a implementação; o segundo pergunta como uma política específica evoluiu e produziu efeitos. Compará-los produtivamente significa manter essa assimetria visível, e não dissolvê-la \parencite{ascenzi2013_implementacao_politicas,mello2017_evolucao_plano}

Também o método diferencia fortemente os textos. No artigo teórico, o procedimento é de revisão analítica e sistematização crítica da literatura, complementado por referências a pesquisas empíricas anteriores que funcionam como sustentação ilustrativa para a proposta de integração conceitual. O ganho desse método está na capacidade de reorganizar o debate do campo e explicitar seus pressupostos. Seu limite é que a proposta alternativa aparece mais como agenda de pesquisa do que como modelo plenamente testado. No estudo sobre o plano de prevenção e controle do desmatamento, o método é documental e longitudinal: os autores comparam publicações científicas sobre o problema, documentos das diferentes fases do plano e avaliações realizadas ao longo de sua implementação. O ganho aqui está na reconstrução histórica e institucional da política, capaz de mostrar mudança, continuidade e aprendizado. O limite correspondente é menor acesso às interações microscópicas da implementação e às interpretações dos agentes de linha de frente \parencite{ascenzi2013_implementacao_politicas,mello2017_evolucao_plano}

No plano conceitual, a comparação revela uma diferença de arquitetura analítica. O artigo sobre perspectivas da implementação organiza o debate em torno de dois grandes modelos: um mais próximo da valorização da hierarquia organizacional, do desenho e da formulação; outro mais atento aos espaços locais de implementação e à atuação das burocracias executoras. Sua intervenção consiste em afirmar que ambos são insuficientes se não incorporarem a dimensão ideacional, isto é, as formas pelas quais os atores interpretam problemas, objetivos e meios de ação. O estudo sobre o plano amazônico, por sua vez, estrutura-se por categorias ligadas à formação de agenda, à coordenação interministerial, ao monitoramento, ao ordenamento territorial, ao comando e controle e à evolução das fases de implementação. Nesse caso, o vocabulário conceitual está mais voltado à dinâmica de uma política pública ambiental concreta e à identificação dos instrumentos mais eficazes. Portanto, um texto produz conceitos para analisar políticas; o outro mobiliza categorias de políticas públicas para interpretar um caso de implementação \parencite{ascenzi2013_implementacao_politicas,mello2017_evolucao_plano}

Os pressupostos teóricos de cada texto também não são idênticos. O artigo de sistematização se aproxima de uma sociologia da ação pública e de uma visão interpretativa da implementação, ao insistir que a ação administrativa não pode ser compreendida apenas por variáveis estruturais, recursos ou desenho institucional. Ele pressupõe que a política é mediada por sentidos, classificações e valores, e que os agentes implementadores não apenas executam, mas interpretam e, nesse movimento, transformam a própria política. O estudo sobre o plano de controle do desmatamento, embora reconheça a complexidade causal do problema e a necessidade de coordenação, opera com pressupostos mais próximos da análise de políticas públicas orientada para capacidades estatais, arranjos institucionais e eficácia de instrumentos. Seu foco recai menos sobre produção de sentido pelos agentes e mais sobre a evolução institucional do plano e sobre a efetividade diferenciada de seus eixos. A tensão entre essas duas posições é fecunda: uma sublinha mediações cognitivas; a outra, capacidades estatais e desenho operacional \parencite{ascenzi2013_implementacao_politicas,mello2017_evolucao_plano}

Há, contudo, uma complementaridade real entre os trabalhos quando se observa o que cada um permite enxergar no outro. O artigo teórico fornece uma lente crítica para reler o estudo do plano amazônico, especialmente ao sugerir que resultados de implementação não dependem apenas de coordenação, monitoramento e coerência institucional, mas também de como os atores interpretam os objetivos da política, os destinatários da intervenção e as formas legítimas de ação estatal. Sob essa chave, o caso do plano poderia ser aprofundado por meio de investigações sobre burocracias implementadoras, conflitos interpretativos entre órgãos, diferenças territoriais de aplicação e formas locais de recepção das medidas. Em sentido inverso, o estudo empírico oferece ao artigo teórico um terreno substantivo onde se pode observar por que variáveis materiais e institucionais continuam centrais: em políticas de grande escala territorial e elevada complexidade federativa, capacidades de monitoramento, fiscalização e coordenação não são meros detalhes administrativos, mas condições decisivas de efetividade \parencite{ascenzi2013_implementacao_politicas,mello2017_evolucao_plano}

A contribuição de cada texto para a disciplina, portanto, se distribui em planos diferentes. O artigo sobre perspectivas analíticas contribui sobretudo para o refinamento conceitual do campo da implementação. Sua principal força é desnaturalizar a oposição entre modelos e demonstrar que a análise da implementação ganha poder explicativo quando articula estrutura organizacional, contexto local e dimensão ideacional. Isso o torna particularmente relevante para pesquisas interessadas em compreender por que políticas formalmente semelhantes produzem trajetórias distintas. Já o estudo do plano de prevenção e controle do desmatamento contribui por mostrar, em uma política ambiental de alta complexidade, como arranjos intersetoriais e instrumentos combinados podem produzir resultados concretos de grande magnitude. Para a disciplina, ele oferece um caso importante de capacidade estatal coordenada, sem ocultar que o êxito foi parcial e mais robusto no controle do que na promoção de alternativas sustentáveis \parencite{ascenzi2013_implementacao_politicas,mello2017_evolucao_plano}

As divergências mais substantivas aparecem quando se observa o que cada texto tende a deixar em segundo plano. No trabalho teórico, a incorporação de ideias, valores e concepções de mundo amplia o horizonte analítico, mas a proposta permanece relativamente aberta quanto aos mecanismos precisos que conectam essa dimensão a recursos, incentivos, normas e estruturas organizacionais. O texto convence sobre a insuficiência das abordagens tradicionais, mas menos sobre como operacionalizar comparativamente sua alternativa. No estudo do plano amazônico, ocorre o movimento inverso: há maior concretude institucional e maior capacidade de demonstrar evolução e resultados, mas menor elaboração sobre os processos interpretativos e sobre as disputas simbólicas que atravessam a implementação em territórios heterogêneos. Em suma, um texto é forte onde o outro é relativamente mais silencioso, e essa assimetria constitui um dos principais achados da comparação \parencite{ascenzi2013_implementacao_politicas,mello2017_evolucao_plano}

Do ponto de vista dos limites comparativos, é importante não forçar uma unidade temática artificial. Os textos pertencem ao mesmo universo disciplinar, mas não foram escritos para responder à mesma pergunta. Um busca reorganizar teoricamente o debate sobre implementação; o outro reconstrói a trajetória de uma política ambiental específica. A comparação, portanto, só é legítima nos eixos transversais realmente sustentados pelo corpus: concepção de implementação, peso de instituições e contextos, tipo de evidência mobilizada, nível de abstração, contribuição analítica e limites. Seria indevido exigir do estudo do plano o mesmo grau de elaboração conceitual sobre ideias e valores que se encontra no artigo teórico, assim como seria inadequado cobrar do texto de revisão a demonstração empírica detalhada de uma política concreta. O ganho comparativo não está em uniformizá-los, mas em mostrar como se situam em pontos diferentes do continuum entre teoria e análise aplicada \parencite{ascenzi2013_implementacao_politicas,mello2017_evolucao_plano}

Sem forçar convergência excessiva, a articulação mais produtiva entre os dois textos consiste em tratá-los como mutuamente esclarecedores. O artigo teórico indica que estudos de implementação ganham densidade quando incorporam as interpretações dos atores; o estudo do plano mostra que a análise de políticas complexas exige atenção rigorosa a capacidades estatais, instrumentos e evolução institucional. Juntos, os textos sugerem que uma análise robusta da implementação precisa combinar, de um lado, sensibilidade às estruturas de coordenação, ao desenho operacional e aos recursos de ação estatal e, de outro, atenção às mediações locais, à discricionariedade e aos sentidos atribuídos pelos agentes à política em execução. Essa articulação não elimina as diferenças entre os trabalhos; ao contrário, preserva-as como fonte de enriquecimento analítico. O corpus, assim, não oferece uma tese única e homogênea, mas um diálogo disciplinar consistente sobre como compreender políticas públicas para além da oposição simplista entre formulação e execução \parencite{ascenzi2013_implementacao_politicas,mello2017_evolucao_plano}



\section{Síntese analítica}
\label{sec:org6f27d0a}

A leitura conjunta do corpus permite sustentar que sua contribuição principal não está em oferecer uma teoria única da implementação, mas em tornar visível que o problema explicativo das políticas públicas se desloca quando se abandona a separação rígida entre desenho e execução. O texto de natureza teórica mostra que a literatura tradicional oscilou entre privilegiar hierarquia, formulação e desenho institucional, de um lado, e valorizar contexto local e discricionariedade burocrática, de outro. Já o estudo sobre o plano amazônico demonstra empiricamente que políticas complexas se estabilizam e produzem efeitos por meio de arranjos institucionais, aprendizado entre fases, coordenação entre órgãos e combinação de instrumentos. Quando lidos em conjunto, os textos sugerem que a implementação é menos uma fase e mais um campo de produção prática da própria política. \parencite{ascenzi2013_implementacao_politicas,mello2017_evolucao_plano}

A interpretação mais produtiva do corpus é a de que ele ilumina uma tensão constitutiva dos estudos de implementação. De um lado, nenhuma política de grande escala pode ser entendida sem capacidades estatais, instrumentos operacionais, coordenação administrativa e mecanismos de monitoramento. O caso do enfrentamento ao desmatamento mostra precisamente isso: a redução do problema dependeu de aparato institucional, integração federativa e intersetorial, além de instrumentos de comando e controle suficientemente robustos para alterar comportamentos e induzir resultados agregados. De outro lado, essas capacidades não esgotam a explicação. O texto teórico insiste que a ação dos implementadores é sempre mediada por esquemas de interpretação, valores e concepções de mundo, de modo que a política não é apenas executada, mas lida, filtrada e redefinida por seus agentes. \parencite{mello2017_evolucao_plano,ascenzi2013_implementacao_politicas}

O ponto decisivo, portanto, não é escolher entre uma explicação institucional e uma explicação interpretativa, mas reconhecer que o corpus evidencia a insuficiência de qualquer modelo unilateral. Se lido isoladamente, o estudo sobre o plano amazônico pode levar a uma ênfase predominante nas capacidades formais do Estado e nos instrumentos mais diretamente observáveis. Se lido isoladamente, o texto sobre perspectivas analíticas pode induzir uma valorização mais abstrata das mediações cognitivas, ainda em registro programático. A leitura articulada corrige os excessos de cada polo: ela mostra que a efetividade estatal exige infraestrutura organizacional e desenho operacional, mas também que esses elementos só ganham realidade concreta quando atravessados por práticas, interpretações e disputas situadas. \parencite{mello2017_evolucao_plano,ascenzi2013_implementacao_politicas}

Essa articulação permite formular uma interpretação própria do corpus: a implementação aparece como uma zona de conversão entre capacidade estatal e construção de sentido. A política pública não se realiza apenas porque foi formalmente desenhada, nem apenas porque agentes locais a reinterpretam; ela se realiza na interação entre arquitetura institucional e mediações práticas. No caso do plano de combate ao desmatamento, isso se expressa na combinação de coordenação central, monitoramento territorial, fiscalização e reconfiguração progressiva das fases do programa. No caso do artigo teórico, essa mesma dinâmica é elevada a problema conceitual, quando os autores mostram que regras, recursos e hierarquia precisam ser pensados junto com ideias, valores e percepções dos implementadores. Em ambos, a política é processo de tradução institucionalmente condicionada. \parencite{mello2017_evolucao_plano,ascenzi2013_implementacao_politicas}

O corpus também permite distinguir entre êxito operacional e transformação substantiva. O estudo sobre o PPCDAm mostra que uma política pode ser bem-sucedida na contenção de um problema por meio de instrumentos de controle e ordenamento, sem que isso signifique ter resolvido as bases estruturais que reproduzem esse problema. Esse achado ganha profundidade quando relido à luz do texto teórico: se a implementação depende de condições materiais e de interpretações dos atores, então políticas baseadas majoritariamente em enforcement tendem a produzir efeitos relevantes, porém potencialmente limitados, quando não conseguem reordenar sentidos, incentivos e práticas de forma mais duradoura. A contribuição do corpus, nesse ponto, é sugerir que controlar um processo não equivale necessariamente a reorganizar o campo social em que ele se produz. \parencite{mello2017_evolucao_plano,ascenzi2013_implementacao_politicas}

Em termos disciplinares, os textos se complementam justamente porque ocupam posições diferentes no continuum entre abstração e concretude. O artigo sobre implementação amplia a sensibilidade teórica do campo e funciona como corretivo contra explicações excessivamente formalistas. O estudo do plano amazônico, por sua vez, impede que a análise da implementação se dissolva em um vocabulário genérico sobre interpretação e contexto, lembrando que políticas públicas dependem de aparatos decisórios, capacidade de coordenação e instrumentos efetivos. A síntese do corpus, assim, não é a fusão indistinta entre duas abordagens, mas a defesa de uma análise em dois níveis: um nível das capacidades e arranjos que tornam a ação possível, e um nível das mediações práticas que definem como essa ação adquire forma concreta. \parencite{ascenzi2013_implementacao_politicas,mello2017_evolucao_plano}



\section{Considerações finais}
\label{sec:orgdf18372}

As análises desenvolvidas ao longo do documento permitem concluir que o principal ganho do corpus está em qualificar o modo de pensar a implementação de políticas públicas. Em vez de tratá-la como desdobramento automático da formulação, os textos mostram que ela envolve coordenação, adaptação, conflito, aprendizado, uso de instrumentos e mediações interpretativas. Essa conclusão é relevante porque recoloca a implementação no centro da análise da ação estatal, não como momento residual, mas como espaço decisivo em que a política adquire densidade institucional e efeitos concretos. \parencite{ascenzi2013_implementacao_politicas,mello2017_evolucao_plano}

A principal contribuição analítica do primeiro texto está em mostrar que o campo de implementação ganha poder explicativo quando incorpora, além das variáveis organizacionais e contextuais, as ideias, os valores e as concepções de mundo dos atores envolvidos. Já o segundo texto contribui ao evidenciar, em um caso substantivo de alta complexidade, que capacidades estatais, coordenação intersetorial e instrumentos adequados podem produzir resultados expressivos, ainda que parciais. Tomados em conjunto, os dois trabalhos reforçam que a disciplina avança mais quando articula reflexão conceitual e exame empírico rigoroso do que quando privilegia unilateralmente um desses polos. \parencite{ascenzi2013_implementacao_politicas,mello2017_evolucao_plano}

Para a disciplina, o corpus oferece ao menos três contribuições duradouras. Primeiro, reafirma a necessidade de superar dicotomias simplificadoras entre formulação e execução. Segundo, demonstra que políticas públicas complexas exigem análise simultânea de desenho institucional, capacidades operacionais e contextos de implementação. Terceiro, sugere que resultados observáveis não podem ser plenamente compreendidos sem atenção às interpretações e práticas dos agentes que materializam a ação estatal. Essas contribuições são particularmente relevantes para a administração pública e para a análise de políticas, pois indicam um programa de pesquisa mais integrado e menos dependente de explicações monocausais. \parencite{ascenzi2013_implementacao_politicas,mello2017_evolucao_plano}

Em síntese, o documento mostra que a heterogeneidade do corpus não é um obstáculo à análise, mas sua própria condição de fecundidade. A diferença entre um texto voltado à reconstrução conceitual e outro dedicado a uma política empírica específica permite enxergar, com mais nitidez, que a implementação é o lugar em que normas, instituições, capacidades, territórios e sentidos se encontram. Esse é o ganho central da leitura realizada: transformar dois textos distintos em um diálogo analítico consistente sobre como compreender, com maior precisão, o funcionamento real das políticas públicas. \parencite{ascenzi2013_implementacao_politicas,mello2017_evolucao_plano}

\printbibliography

\clearpage
\section{Mapa mental dos textos analisados}
\label{sec:orgee1a65e}
\vspace{0.8em}
\begingroup
\centering
\includegraphics[width=\textwidth,height=0.78\textheight,keepaspectratio]{images/mapa_mental.png}
\par
\endgroup
\clearpage
\end{document}